%! Logarithmic Function Manipulation — Unit 5, Lesson 5.4 — Group Activity
\documentclass[10pt]{article}
\usepackage{precalculus-article}
\usepackage{precalculus-boxes}
\newcommand{\TallMath}[1]{$\displaystyle #1\rule[-1.4em]{0pt}{3.2em}$}

\begin{document}

\pageheader{Unit 5, Lesson 5.4}{Group Activity}
\namepartnerperiod

\vspace{0.1in}

\begin{scenariobox}[Signal Power Analysis]{sky}
A radio engineer uses the formula $\text{dB} = 10 \cdot \log_{10}(P/P_0)$ to measure signal
power levels. Your group will simplify and rewrite logarithmic expressions to help her
analyze signals with different power levels.
\end{scenariobox}

\vspace{0.1in}

\begin{tcolorbox}[colback=white, colframe=black!40,
    title=\textbf{Tier R --- Remediate},
    fonttitle=\bfseries, arc=2mm, left=3mm, right=3mm, top=2mm, bottom=2mm]

\textbf{Apply the product and power properties to numerical expressions.}

\begin{enumerate}[leftmargin=1.8em, itemsep=10pt]

\item Rewrite each expression as a sum using the product property.

\vspace{4pt}
\begin{enumerate}[leftmargin=1.5em, itemsep=8pt, label=\alph*.]
  \item $\log_2(4 \cdot 8) = \log_2(\blank{1.5cm}) + \log_2(\blank{1.5cm}) = \blank{1cm} + \blank{1cm} = \blank{1.5cm}$
  \item $\log_3(9 \cdot 27) = \log_3(\blank{1.5cm}) + \log_3(\blank{1.5cm}) = \blank{1cm} + \blank{1cm} = \blank{1.5cm}$
  \item $\log_5(25 \cdot 125) = \log_5(\blank{1.5cm}) + \log_5(\blank{1.5cm}) = \blank{1cm} + \blank{1cm} = \blank{1.5cm}$
\end{enumerate}

\vspace{4pt}
\item Apply the power property to rewrite each expression.

\vspace{4pt}
\begin{enumerate}[leftmargin=1.5em, itemsep=8pt, label=\alph*.]
  \item $\log_2(4^2) = \blank{1cm} \cdot \log_2(\blank{1.5cm}) = \blank{1cm} \cdot \blank{1cm} = \blank{1.5cm}$
  \item $\log_3(27^3) = \blank{1cm} \cdot \log_3(\blank{1.5cm}) = \blank{1cm} \cdot \blank{1cm} = \blank{1.5cm}$
\end{enumerate}

\end{enumerate}
\end{tcolorbox}

\vspace{0.1in}

\begin{tcolorbox}[colback=white, colframe=black!40,
    title=\textbf{Tier A --- Approaching Proficiency},
    fonttitle=\bfseries, arc=2mm, left=3mm, right=3mm, top=2mm, bottom=2mm]

\textbf{Rewrite algebraic expressions and apply change of base.}

\begin{enumerate}[leftmargin=1.8em, itemsep=10pt]

\item Rewrite $\log_2(32x)$ in the form $a + \log_2(x)$. What is the value of $a$?

\vspace{4pt}
$\log_2(32x) = \log_2(\blank{2cm}) + \log_2(x) = \blank{1.5cm} + \log_2(x)$

\vspace{2pt}
$a = $ \blank{2cm}. This represents a vertical \blank{3cm} of the graph of $y = \log_2(x)$.

\item Rewrite $\log_3(x^4)$ as a vertical dilation of $\log_3(x)$.

\vspace{4pt}
$\log_3(x^4) = $ \blank{4cm}. The dilation factor is \blank{2cm}.

\item Use change of base to evaluate $\log_5(17)$ using common logarithms ($\log_{10}$).

\vspace{4pt}
$\log_5(17) = \dfrac{\log(\blank{2cm})}{\log(\blank{2cm})} \approx $ \blank{2.5cm}

\item Show that $\log_2(x) = \dfrac{\ln(x)}{\ln(2)}$ using the change-of-base formula.

\vspace{4pt}
\writelines{3}

\end{enumerate}
\end{tcolorbox}

\vspace{0.1in}

\begin{tcolorbox}[colback=white, colframe=black!40,
    title=\textbf{Tier E --- Extension},
    fonttitle=\bfseries, arc=2mm, left=3mm, right=3mm, top=2mm, bottom=2mm]

\textbf{Prove and justify log property connections to graph transformations.}

\begin{enumerate}[leftmargin=1.8em, itemsep=10pt]

\item Prove that $f(x) = \log_3(9x) = g(x) + 2$ where $g(x) = \log_3(x)$.
      What geometric transformation does this represent? Justify your reasoning.

\vspace{4pt}
\writelines{4}

\item Use the change-of-base formula to show that $\log_b(x) = \dfrac{\log_{10}(x)}{\log_{10}(b)}$.
      Explain why this means every logarithmic function is a vertical dilation of $\log_{10}(x)$.

\vspace{4pt}
\writelines{4}

\item Simplify $\log_2(8x^3)$ completely using both the product and power properties. Show
      all steps and name each property as you use it.

\vspace{4pt}
\writelines{4}

\end{enumerate}
\end{tcolorbox}

\end{document}
